\documentclass[12pt]{article}

\usepackage[spanish]{babel}
\usepackage[utf8]{inputenc}
\usepackage{amsmath}
\usepackage{amssymb}
\usepackage{xcolor}
\usepackage{listings}

\lstset{
language=R,
basicstyle=\ttfamily\small,
numbers=left,
numberstyle=\tiny,
commentstyle=\color{green!50!black},
keywordstyle=\color{blue},
stringstyle=\color{red},
breaklines=true,
frame=single
}

\begin{document}

%================ PORTADA ================
\begin{center}

{\LARGE \textbf{EJERCICIOS DE ANUALIDADES}}\\[0.5cm]

{\large Estadística Actuarial}\\[1cm]

\textbf{Estudiante:} Kely Rosio Sanchez Champi\\[0.3cm]

\end{center}

\vspace{1cm}

%================ CONTENIDO ================


\section*{Ejercicio 1: Valor Actual de una Anualidad Vencida}

Una persona adquiere una motocicleta mediante un financiamiento que consiste en realizar pagos mensuales de S/1\,800 al final de cada mes durante 4 años. Si la entidad financiera aplica una tasa nominal anual del 12\% con capitalización mensual, determinar el valor actual de la deuda.

\subsection*{Datos}

\begin{itemize}
\item \textbf{Valor actual (A):} ?
\item \textbf{Renta o pago mensual (R):} S/1\,800
\item \textbf{Tiempo de financiamiento:} 4 años
\[
n=(4)(12)=48\text{ meses}
\]
\item \textbf{Tasa nominal anual:} 12\%
\item \textbf{Capitalización:} Mensual
\item \textbf{Tasa efectiva mensual:}
\[
i=\frac{0.12}{12}=0.01
\]
\end{itemize}

\subsection*{Fórmula}

Para una anualidad vencida:

\[
A=R\left[\frac{1-(1+i)^{-n}}{i}\right]
\]

\subsection*{Desarrollo}

Sustituyendo los valores:

\[
A=1800
\left[
\frac{1-(1+0.01)^{-48}}
{0.01}
\right]
\]

\[
A=1800
\left[
\frac{1-0.62065}
{0.01}
\right]
\]

\[
A=1800(37.9354)
\]

\[
A=68\,283.72
\]

\subsection*{Respuesta}

Por lo tanto, el valor actual de la deuda es:

\[
\boxed{A=S/68\,283.72}
\]

\subsection*{Solución en R}

\begin{lstlisting}
#-------------------------------------------------
# EJERCICIO 1
# Valor actual de una anualidad vencida
#-------------------------------------------------

# Pago mensual
R <- 1800

# Numero de pagos
n <- 48

# Tasa mensual
i <- 0.12/12

# Valor actual
VA <- R*((1-(1+i)^(-n))/i)

# Resultado
print(paste("Valor actual:", round(VA,2)))


\end{lstlisting}

\section*{Ejercicio 2: Valor Futuro de una Anualidad Vencida}

Una trabajadora deposita S/.900 al final de cada mes durante 5 años en una cuenta de ahorros que paga una tasa nominal anual del 9\% con capitalización mensual. Determinar el monto acumulado.

\subsection*{Datos}

\begin{itemize}
\item Valor futuro: \(S=?\)
\item Depósito mensual: \(R=S/.900\)
\item Tiempo: 5 años
\[
n=5(12)=60
\]
\item Tasa nominal anual: 9\%
\[
i=\frac{0.09}{12}=0.0075
\]
\end{itemize}

\subsection*{Fórmula}

\[
S=R\left[\frac{(1+i)^n-1}{i}\right]
\]

\subsection*{Desarrollo}

\[
S=900\left[\frac{(1.0075)^{60}-1}{0.0075}\right]
\]

\[
S=68\,354.11
\]

\subsection*{Respuesta}

\[
\boxed{S=S/.68\,354.11}
\]

\subsection*{Código en R}

\begin{lstlisting}
# Ejercicio 2

R <- 900
n <- 60
i <- 0.09/12

VF <- R*(((1+i)^n-1)/i)

round(VF,2)
\end{lstlisting}

\newpage

\section*{Ejercicio 2: Cálculo de la Renta en una Anualidad Vencida}

Una persona recibe un préstamo cuyo valor actual es de S/50\,000. El financiamiento se pagará mediante cuotas mensuales iguales durante 3 años. Si la tasa nominal anual es 18\% con capitalización mensual, determinar el valor de la cuota mensual.

\subsection*{Datos}

\begin{itemize}
\item \textbf{Valor actual (A):} S/50\,000
\item \textbf{Renta (R):} ?
\item \textbf{Tiempo:}
\[
n = 3 \times 12 = 36 \text{ meses}
\]
\item \textbf{Tasa nominal anual:} 18\%
\item \textbf{Tasa mensual:}
\[
i = \frac{0.18}{12} = 0.015
\]
\end{itemize}

\subsection*{Fórmula}

\[
A = R\left[\frac{1-(1+i)^{-n}}{i}\right]
\]

Despejando R:

\[
R = \frac{A}{\left[\frac{1-(1+i)^{-n}}{i}\right]}
\]

\subsection*{Desarrollo}

\[
R = \frac{50\,000}{\left[\frac{1-(1.015)^{-36}}{0.015}\right]}
\]

\[
(1.015)^{-36} \approx 0.5962
\]

\[
R = \frac{50\,000}{\left[\frac{1-0.5962}{0.015}\right]}
\]

\[
R = \frac{50\,000}{26.92}
\]

\[
R \approx 1\,857.73
\]

\subsection*{Respuesta}

\[
\boxed{R \approx S/1\,857.73}
\]

\subsection*{Solución en R (desarrollo detallado)}

\begin{lstlisting}
#-------------------------------------------------
# EJERCICIO 2
# Cálculo de la renta en una anualidad vencida
#-------------------------------------------------

# Datos financieros
A <- 50000          # Valor actual del prestamo
n <- 36             # Numero de meses
i <- 0.18/12        # Tasa mensual

# Paso 1: Factor de valor actual de una anualidad
factor_VA <- (1 - (1+i)^(-n)) / i

# Mostrar factor intermedio
print(paste("Factor de valor actual:", round(factor_VA,4)))

# Paso 2: Calculo de la renta
R <- A / factor_VA

# Resultado final
print(paste("Cuota mensual R:", round(R,2)))

# Paso 3: Verificacion (reconstruccion del valor actual)
A_verificado <- R * factor_VA
print(paste("Verificacion VA:", round(A_verificado,2)))

\end{lstlisting}

\section*{Ejercicio 3: Cálculo del Tiempo en una Anualidad Vencida}

Un préstamo de S/30\,000 se cancela mediante pagos mensuales de S/1\,200 al final de cada mes. Si la tasa nominal anual es 12\% capitalizable mensualmente, determinar el número de meses necesarios para cancelar la deuda.

\subsection*{Datos}

\begin{itemize}
\item \textbf{Valor actual (A):} S/30\,000
\item \textbf{Renta (R):} S/1\,200
\item \textbf{Tiempo (n):} ?
\item \textbf{Tasa nominal anual:} 12\%
\item \textbf{Tasa mensual:}
\[
i = \frac{0.12}{12} = 0.01
\]
\end{itemize}

\subsection*{Fórmula}

\[
A = R\left[\frac{1-(1+i)^{-n}}{i}\right]
\]

Despejando \(n\):

\[
(1+i)^{-n} = 1 - \frac{Ai}{R}
\]

\[
-n \ln(1+i) = \ln\left(1 - \frac{Ai}{R}\right)
\]

\[
n = \frac{-\ln\left(1 - \frac{Ai}{R}\right)}{\ln(1+i)}
\]

\subsection*{Desarrollo}

\[
n = \frac{-\ln\left(1 - \frac{30\,000(0.01)}{1\,200}\right)}{\ln(1.01)}
\]

\[
n = \frac{-\ln(1 - 0.25)}{\ln(1.01)}
\]

\[
n = \frac{-\ln(0.75)}{\ln(1.01)}
\]

\[
n \approx \frac{0.28768}{0.00995}
\]

\[
n \approx 28.9 \approx 29 \text{ meses}
\]

\subsection*{Respuesta}

\[
\boxed{n \approx 29 \text{ meses}}
\]

\subsection*{Solución en R (desarrollo detallado)}

\begin{lstlisting}
#-------------------------------------------------
# EJERCICIO 3
# Cálculo del número de meses en una anualidad vencida
#-------------------------------------------------

# Datos financieros
A <- 30000          # Valor actual
R <- 1200           # Pago mensual
i <- 0.12/12        # Tasa mensual

# Paso 1: Proporcion del valor presente
proporcion <- 1 - (A*i)/R

print(paste("Proporcion calculada:", round(proporcion,4)))

# Validacion de consistencia financiera
if (proporcion <= 0) {
  print("Error: la renta no cubre el valor del prestamo con esta tasa.")
} else {

  # Paso 2: Calculo del numero de periodos
  n <- -log(proporcion) / log(1+i)

  # Resultado
  print(paste("Numero de meses (n):", round(n,2)))

  # Paso 3: Redondeo financiero (meses completos)
  n_redondeado <- ceiling(n)
  print(paste("Meses redondeados:", n_redondeado))

  # Paso 4: Verificacion del valor actual aproximado
  A_verificado <- R * ((1 - (1+i)^(-n)) / i)
  print(paste("Verificacion VA:", round(A_verificado,2)))
}
\end{lstlisting}
\section*{Ejercicio 4 y 5: Anualidades Anticipadas}

\section*{Ejercicio 4: Valor Actual de una Anualidad Anticipada}

Una persona realiza pagos de S/2\,000 al inicio de cada mes durante 5 años para cancelar una deuda. Si la tasa nominal anual es 12\% capitalizable mensualmente, determinar el valor actual de la deuda.

\subsection*{Datos}

\begin{itemize}
\item \textbf{Renta (R):} S/2\,000
\item \textbf{Tiempo:}
\[
n = 5 \times 12 = 60
\]
\item \textbf{Tasa mensual:}
\[
i = 0.12/12 = 0.01
\]
\end{itemize}

\subsection*{Fórmula (Anualidad anticipada)}

\[
A = R \left[\frac{1-(1+i)^{-n}}{i}\right](1+i)
\]

\subsection*{Desarrollo}

\[
A = 2000 \left[\frac{1-(1.01)^{-60}}{0.01}\right](1.01)
\]

\[
(1.01)^{-60} \approx 0.5488
\]

\[
A = 2000 \left[\frac{1-0.5488}{0.01}\right](1.01)
\]

\[
A = 2000(45.12)(1.01)
\]

\[
A \approx 91\,241.00
\]

\[
\boxed{A \approx S/91\,241.00}
\]

\section*{Ejercicio 5: Cálculo de la Renta en Anualidad Anticipada}

Un préstamo de S/80\,000 debe cancelarse en 4 años mediante pagos mensuales anticipados. Si la tasa nominal anual es 15\% capitalizable mensualmente, determinar la cuota mensual.

\subsection*{Datos}

\begin{itemize}
\item \textbf{Valor actual (A):} S/80\,000
\item \textbf{Tiempo:}
\[
n = 4 \times 12 = 48
\]
\item \textbf{Tasa mensual:}
\[
i = 0.15/12 = 0.0125
\]
\end{itemize}

\subsection*{Fórmula}

\[
A = R \left[\frac{1-(1+i)^{-n}}{i}\right](1+i)
\]

Despejando \(R\):

\[
R = \frac{A}{\left[\frac{1-(1+i)^{-n}}{i}\right](1+i)}
\]

\subsection*{Desarrollo}

\[
R = \frac{80\,000}{\left[\frac{1-(1.0125)^{-48}}{0.0125}\right](1.0125)}
\]

\[
(1.0125)^{-48} \approx 0.5507
\]

\[
R = \frac{80\,000}{(35.94)(1.0125)}
\]

\[
R \approx 2\,197.60
\]

\[
\boxed{R \approx S/2\,197.60}
\]

\subsection*{Solución en R}

\begin{lstlisting}
#=================================================
# EJERCICIO 4 Y 5 - ANUALIDADES ANTICIPADAS
#=================================================

#-------------------------
# EJERCICIO 4
#-------------------------

R1 <- 2000
n1 <- 60
i1 <- 0.12/12

factor1 <- (1 - (1+i1)^(-n1)) / i1
A1 <- R1 * factor1 * (1+i1)

print(paste("VA Ejercicio 4:", round(A1,2)))


#-------------------------
# EJERCICIO 5
#-------------------------

A2 <- 80000
n2 <- 48
i2 <- 0.15/12

factor2 <- (1 - (1+i2)^(-n2)) / i2
R2 <- A2 / (factor2 * (1+i2))

print(paste("R Ejercicio 5:", round(R2,2)))
\end{lstlisting}
\end{document}